\documentclass[11pt]{article}

% ======================================================================
%  Documentation for cubesat_sim - the CubeSat double-Fourier
%  interferometer instrument simulator. Explains the code structure,
%  every module's job, the physics each one implements, the design
%  decisions, and the open questions the build surfaced.
%  Paste into Overleaf and compile with pdfLaTeX.
% ======================================================================

\usepackage[a4paper,margin=1in,headsep=14pt]{geometry}

% ---- type ----------------------------------------------------------------
\usepackage[T1]{fontenc}
\usepackage{mathtools}
\usepackage{libertine}
\usepackage{libertinust1math}
\usepackage{microtype}

\usepackage{siunitx}
\usepackage{xcolor}
\usepackage{booktabs}
\usepackage{array}
\usepackage{graphicx}

% ---- palette (matches the working-notes theme) ----------------------------
\definecolor{plumdark}{HTML}{482661}
\definecolor{plummid}{HTML}{7A4FB0}
\definecolor{plumaccent}{HTML}{963A8A}
\definecolor{plumrule}{HTML}{D6C9EC}
\definecolor{plumtint}{HTML}{F4F0FA}
\definecolor{plumtint2}{HTML}{EDE6F7}
\definecolor{plumtint3}{HTML}{F6EFF4}
\definecolor{ink}{HTML}{20131F}

\AtBeginDocument{\color{ink}}

% ---- boxes -----------------------------------------------------------------
\usepackage{tcolorbox}
\tcbuselibrary{skins,breakable}

\newtcolorbox{physicsbox}[1][]{
  enhanced, breakable, sharp corners, boxrule=0pt, leftrule=2.6pt,
  colframe=plummid, colback=plumtint2, coltitle=plummid,
  fonttitle=\sffamily\bfseries, left=12pt, right=12pt, top=6pt, bottom=9pt,
  before skip=9pt, after skip=9pt,
  title={\textsc{The physics this module implements}}, #1}

\newtcolorbox{decisionbox}[1][]{
  enhanced, breakable, sharp corners, boxrule=0pt, leftrule=2.6pt,
  colframe=plumdark, colback=plumtint, coltitle=plumdark,
  fonttitle=\sffamily\bfseries, left=12pt, right=12pt, top=6pt, bottom=9pt,
  before skip=9pt, after skip=9pt,
  title={\textsc{Design decision}}, #1}

\newtcolorbox{flagbox}[1][]{
  enhanced, breakable, sharp corners, boxrule=0pt, leftrule=2.6pt,
  colframe=plumaccent, colback=plumtint3, coltitle=plumaccent,
  fonttitle=\sffamily\bfseries, left=12pt, right=12pt, top=6pt, bottom=9pt,
  before skip=9pt, after skip=9pt,
  title={\textsc{Flag --- open question / inconsistency found}}, #1}

% ---- lists / headings -------------------------------------------------------
\usepackage{enumitem}
\setlist[itemize,1]{label={\textcolor{plummid}{\footnotesize$\blacksquare$}}, leftmargin=1.5em}
\setlist[enumerate,1]{label={\textcolor{plummid}{\sffamily\bfseries\arabic*.}}, leftmargin=1.7em}

\usepackage{titlesec}
\titleformat{\section}
  {\sffamily\LARGE\bfseries\color{plumdark}}{}{0pt}
  {\textcolor{plummid}{\rule[0.12ex]{0.55em}{0.55em}}\hspace{0.6em}}
  [\vspace{4pt}{\color{plumrule}\titlerule[1.2pt]}]
\titleformat{\subsection}
  {\sffamily\large\bfseries\color{plumdark}}
  {\textcolor{plummid}{\rule[-0.18ex]{2.6pt}{1.05em}}\hspace{0.55em}\textcolor{plummid}{\thesubsection}}
  {0.6em}{}
\titleformat{\paragraph}[runin]{\sffamily\bfseries\color{plumdark}}{}{0pt}{}
\titlespacing*{\paragraph}{0pt}{1.4ex plus .4ex}{0.7em}

\usepackage{fancyhdr}
\pagestyle{fancy}
\fancyhf{}
\fancyhead[L]{\footnotesize\sffamily\color{plumdark}\nouppercase{\leftmark}}
\fancyhead[R]{\footnotesize\sffamily\color{plummid}\thepage}
\renewcommand{\headrule}{{\color{plumrule}\rule{\headwidth}{0.7pt}}}

\usepackage[colorlinks=true, linkcolor=plumdark, citecolor=plummid,
            urlcolor=plummid]{hyperref}

% ---- code style -------------------------------------------------------------
\newcommand{\code}[1]{\texttt{\small #1}}
\newcommand{\file}[1]{\textcolor{plumdark}{\texttt{\small #1}}}

\begin{document}

% ---- title block -------------------------------------------------------------
\begin{center}
{\color{plummid}\rule{\textwidth}{2.4pt}}\\[2.5pt]
{\color{plumrule}\rule{\textwidth}{0.6pt}}\\[16pt]
{\sffamily\bfseries\color{plumdark}\Huge cubesat\_sim}\\[10pt]
{\sffamily\Large\color{plummid} An End-to-End Instrument Simulator for the
 CubeSat\\[3pt] Double-Fourier Spatio-Spectral Interferometer}\\[8pt]
{\itshape\normalsize\color{ink} Architecture, physics, design decisions and
 open questions}\\[14pt]
{\color{plumrule}\rule{\textwidth}{0.6pt}}\\[2.5pt]
{\color{plummid}\rule{\textwidth}{2.4pt}}\\[9pt]
{\sffamily\small\color{plumdark!75} Savini group CubeSat interferometer
 project \;\textcolor{plummid}{\textbullet}\; UCL \;\textcolor{plummid}{\textbullet}\; \today}
\end{center}
\vspace{1.0em}

\begin{tcolorbox}[enhanced, breakable, sharp corners, boxrule=0pt, leftrule=2.6pt,
  colframe=plummid, colback=plumtint2, left=13pt, right=13pt, top=8pt, bottom=9pt]
{\sffamily\bfseries\color{plumdark}\textsc{Summary}}\par\smallskip\small\itshape
\code{cubesat\_sim} is a Python~3 package that simulates the CubeSat
double-Fourier interferometer end to end: a known scene is pushed through a
software model of the instrument to produce realistic detector interferograms,
which are then reduced back to the measured quantities and checked against the
known input. For the confirmed single-fixed-baseline mission that quantity is
the visibility spectrum: detection and one-axis sizing of a feature, not an
image. A multi-baseline mode (baseline rotation and contraction) that
reconstructs source positions and per-source spectra, and compares them with
the input scene, is included as a future/non-mission configuration. It is a
ground-up re-architecture of PyFIInS (the Python port of
R.~Juanola-Parramon's FIInS simulator, FISICA project) for a two-element,
single-baseline, warm 6U platform. The physics follows Grainger et al.\ 2012
(arXiv:1203.2144); every instrument number is traced to a cell of
\code{Cubesat\_interferometer\_v0.2.xlsx} or to the optical layout slides.
Writing and validating the code surfaced six concrete design questions,
including one genuine finding: the spreadsheet's single-sided scan budget
cannot support complex visibility recovery at the design resolution.
\end{tcolorbox}

\tableofcontents
\vspace{0.6em}
\noindent{\color{plumrule}\rule{\textwidth}{0.6pt}}
\vspace{0.4em}

% ======================================================================
\section{Purpose and philosophy}
% ======================================================================

The simulator exists to close one loop. A measurement technique is
trustworthy only if, when handed a known input, it returns that input. The
Grainger laboratory bench closed this loop in hardware: a lamp behind an
aperture of known shape was measured by the interferometer and the recovered
quantities compared with the computable truth. The FIInS simulator closed the
same loop in software for the FISICA flagship concept. \code{cubesat\_sim}
closes it for the CubeSat: \emph{known scene in $\rightarrow$ simulated
detector data $\rightarrow$ recovered quantity out $\rightarrow$ compare}. For
the confirmed single-fixed-baseline mission the recovered quantity is the
visibility spectrum (detection and one-axis sizing of a feature, not an
image); the future multi-baseline mode instead recovers source positions and
per-source spectra. Where input and output match, the design works; where
they diverge, an instrument limitation has been located and quantified.

Two modes of the loop matter, and the code keeps them strictly separate:

\begin{itemize}
\item \textbf{Noiseless mode} (\code{seed=None}) validates the
\emph{architecture}. With no noise, any input/output discrepancy is a code or
physics error, never statistics. The shipped test suite requires the noiseless
loop to close to numerical precision (it does: positions to
$<10^{-3}$~arcsec, spectra to $<10^{-6}$ fractional error).
\item \textbf{Noisy mode} estimates \emph{performance}: thermal background,
photon noise and detector noise are injected at the radiometric model's
levels, and the degradation of the recovered scene measures what the real
instrument would deliver.
\end{itemize}

This separation is what makes the simulator a design tool rather than a demo:
parameters become knobs, and the validated loop converts each knob setting
into a science outcome.

% ======================================================================
\section{What is being modelled}
% ======================================================================

The instrument, as defined by the v0.2 radiometric spreadsheet and the
optical layout slides: a 6U CubeSat payload carrying two 2-inch
(\SI{50.8}{\milli\metre}) collecting telescopes, each feeding a beam-reducer
chain into a central Michelson-type combiner; one arm of the combiner carries
the scanning roof-mirror pair that sweeps the optical path difference (OPD);
a single detector records one power sample per OPD step. Band
\SIrange{8}{12}{\micro\metre} at spectral resolution $R=100$; the whole
instrument is warm (\SI{293}{\kelvin}) --- significant because Wien's law puts
the peak of \SI{293}{\kelvin} thermal emission at \SI{9.9}{\micro\metre},
squarely in band, so the instrument glows at itself.

Spatial information enters through the baseline $b$ between the telescopes.
The confirmed design is a SINGLE FIXED baseline of ${\sim}\SI{0.4}{\metre}$
(detection + one-axis sizing). The spreadsheet's \SIrange{0.3}{0.5}{\metre}
with contraction is a future multi-baseline option, and the rigid bench is
\SI{125}{\milli\metre} (see Flag~1); spectral information enters through the
scanned OPD $\Delta$. In that future multi-baseline mode the uv plane is
filled by spinning the spacecraft about the line of sight and stepping the
baseline length between half-turns; the confirmed single-baseline instrument
does neither. Because the
instrument is a two-element single-baseline system, its realistic data product
is \emph{source separation plus spectroscopy} --- the double-slit regime of
Grainger Eqs.~11--13 --- not wide-field image synthesis. The code's science
pipeline is built around exactly that product.

\begin{decisionbox}[title={\textsc{What the code takes from the optical
layout slides}}]
Five facts of the model come from the layout rather than the spreadsheet:
(i) the \SI{125}{\milli\metre} bench separation behind Flag~1; (ii) the
roof-pair delay line, fixing mechanical travel $=\Delta/2$ (the
``$\tfrac12\,\delta$OPD'' annotation); (iii) the single detector (+TEC);
(iv) the asymmetric extra fold in one routing, behind the per-arm
throughputs and the parity question; and (v) the path budgets of slides
3--4, which \emph{balance by design}: summing the C and D rows gives
$10{+}10{+}5{+}4 = 29$ on the right routing and
$10{+}5{+}3{+}5{+}6{+}\delta = 29{+}\delta$ on the left --- the extra fold
is path-length compensated and ZPD sits at the stage's nominal zero. The
design ZPD offset is therefore zero; the \emph{residual} offset
(manufacturing, thermal drift) is the real instrument error, carried as
\code{FTSScan.zpd\_offset\_cm}: the forward model shifts the white-light
fringe by it, and the reduction removes it only if told its value, so an
uncalibrated ZPD error and its phase-ramp signature can be studied
directly. What the layout offers but the code does not yet use: the
specific OAP beam-reducer prescriptions (radii/focal lengths), and beam
walk on the roof mirror during the scan.
\end{decisionbox}

\begin{decisionbox}[title={\textsc{What the code takes from the radiometric
model}}]
The spreadsheet is the codebase's backbone: 81 cited \code{[RM]} sheet/cell
references across the source. Every number in \code{Band}, \code{Telescope},
\code{ColdOptics}, \code{Detector}, \code{FTSScan}, \code{Geometry} and
\code{Environment} traces to it; its three background loadings are
reproduced at ratio 1.00 (test-enforced to 5\%); its uv strategy (spin
rates, contraction, radial density) is the schedule; its data-rate sheet
(16 bit $\times$ \SI{666.7}{\hertz}) gives the science-volume estimate;
and its \emph{empty} rows are implemented rather than skipped --- the
sensitivity chain (NEP$_\mathrm{tot}$, noise per channel, 5$\sigma$ MDFD),
the parked ``visibility loss due to intensity mismatch'' term, and a
saturation-power hook ([PH], warns when a sample clips). Two of its numbers
are deliberately \emph{rejected}, with reasons: the photon-NEP cells
(internally inconsistent with its own loadings, Flag~5) and the
single-sided scan budget (cannot carry complex visibilities, Flag~4). Its
mass/power tables (FISICA strawman leftovers), its \SI{1800}{\kelvin}
lab-testbed source scene, and the zeroed channels B2--B4 are knowingly
unused.
\end{decisionbox}

% ======================================================================
\section{Package structure}
% ======================================================================

\begin{verbatim}
cubesat_codebase/
|-- run_simulation.py        end-to-end demo: full multi-baseline uv pass + imaging (future/non-mission)
|-- run_cases.py             single-baseline demo: the three targets
|-- run_crater.py            single-baseline lunar-crater detection study
|-- requirements.txt         numpy, scipy, matplotlib, pytest
|-- README.md                quick start + open-questions list
|-- outputs/                 figures (full pass; outputs/cases/, crater/)
|-- tests/
|   |-- test_physics.py      hand-checkable physics tests (full pipeline)
|   |-- test_singlebaseline.py  single-baseline + scene-case tests
|   '-- test_extended.py     extended-scene loader (vs discrete limit)
'-- cubesat_sim/
    |-- constants.py         physical constants, Planck, unit conversions
    |-- parameters.py        ALL instrument numbers, source-traced
    |-- fts.py               wavenumber + OPD sampling grids
    |-- sky.py               scene model (compact sources, spectra)
    |-- scenes.py            three runnable target cases (the building blocks)
    |-- extended.py          pixel-map scenes (Moon patch, crater)
    |-- beams.py             Airy beam; parity-mismatch envelope
    |-- uvcoverage.py        spin + contraction baseline schedule
    |-- radiometry.py        background, photon/detector NEP, sensitivity
    |-- forward.py           scene -> visibilities -> interferograms
    |-- reduction.py         interferograms -> visibility spectra
    |-- imaging.py           dirty map; positions + spectra extraction
    |-- singlebaseline.py    one baseline: interferogram <-> spectrum
    '-- validate.py          input-vs-recovered metrics and figures
\end{verbatim}

The data flow is a single chain; each module consumes the previous module's
output through a small, explicit interface (a dataclass), so any stage can be
modified or replaced without touching the others --- the modularity principle
inherited from FIInS:

\begin{center}\sffamily\small
parameters $\rightarrow$ grids $\rightarrow$ uv schedule $\rightarrow$ scene
$\rightarrow$ \textbf{forward} $\rightarrow$ FITS-like scan list
$\rightarrow$ \textbf{reduction} $\rightarrow$ visibility spectra
$\rightarrow$ \textbf{imaging} $\rightarrow$ positions + spectra
$\rightarrow$ \textbf{validate}
\end{center}

% ======================================================================
\section{Module by module}
% ======================================================================

% ----------------------------------------------------------------------
\subsection{\file{constants.py} --- units and the Planck function}

Defines the package conventions: wavenumber $\sigma$ in \si{\per\centi\metre}
(so a wavenumber $\sigma$ produces a fringe $\cos(2\pi\sigma\Delta)$ with OPD
$\Delta$ in \si{\centi\metre}, with no stray conversion factors); angles in
radians internally, arcsec at interfaces; spectral radiance per
\si{\per\centi\metre}. Provides the Planck function in those units,
\begin{equation}
B_\sigma(T) \;=\; \frac{2 h \nu^3}{c^2}\,
\frac{1}{e^{h\nu/kT}-1}\;\cdot\;100\,c,
\qquad \nu = 100\,c\,\sigma,
\end{equation}
and the Jansky conversion $1\,\mathrm{Jy} = 10^{-26}\cdot 100\,c$
\,\si{\watt\per\square\metre} per \si{\per\centi\metre}.

% ----------------------------------------------------------------------
\subsection{\file{parameters.py} --- every number, source-traced}

All instrument numbers live here, in flat dataclasses (\code{Band},
\code{Telescope}, \code{ColdOptics}, \code{Detector}, \code{FTSScan},
\code{Geometry}, \code{Environment}), each value carrying a comment with its
source: \code{[RM]} sheet/cell of the radiometric spreadsheet, \code{[OPT]}
slide of the optical layout, or \code{[PH]} placeholder. The file contains no
physics --- only numbers and trivially derived quantities --- so that a
parameter change can never silently alter an equation.

Key traced values: $d_\mathrm{dish}=\SI{50.8}{\milli\metre}$, telescope
efficiency $0.98^2 = 0.9801$, cold-optics throughput
$0.98^5 \times 0.49 \times 0.98 \times 0.95 = 0.4124$ (the $0.49$ is the beam
combiner and is included here exactly as the spreadsheet does, so it is never
double-counted), overall optical efficiency $0.404$, detector absorption
$0.9$, cold-optics equivalent emissivity $0.0811$, telescope equivalent
emissivity $0.0199$.

\begin{decisionbox}
\textbf{The asymmetric extra fold mirror} ([OPT] slides 3--4: path budgets
$10{+}10$ and $3{+}5{+}6{+}\delta$ on one routing versus $10{+}5$ and $5{+}4$
on the other) is modelled as per-arm throughputs $t_1 = 1$,
$t_2 = \rho^{N_\mathrm{extra}} = 0.98$. The resulting fringe-contrast loss is
$2\sqrt{t_1 t_2}/(t_1+t_2) = 0.99995$ --- four parts in $10^5$, i.e.\
\emph{negligible}, which is itself the quantitative answer to ``does the
extra mirror's throughput matter?'' (No. Its parity might; see
\file{beams.py}.) The spreadsheet keeps one symmetric throughput and lists
``visibility loss due to intensity mismatch'' under \emph{things to add}; the
code therefore goes one step beyond it here.
\end{decisionbox}

% ----------------------------------------------------------------------
\subsection{\file{fts.py} --- the sampling grids}

\begin{physicsbox}
Standard FTS sampling theory (identical logic to PyFIInS \code{fts.py}).
Channel width from the resolution requirement:
$\delta\sigma = \sigma_c/R = \SI{10}{\per\centi\metre}$. Maximum OPD from the
Fourier reciprocal relation: $\Delta_{\max} = 1/(2\,\delta\sigma) =
\SI{0.05}{\centi\metre}$. OPD step from Nyquist at the band top, with the
spreadsheet's $\times 2$ oversampling: $\delta\Delta =
1/(2\sigma_{\max}\cdot 2) = \SI{2}{\micro\metre}$. These reproduce
[RM]~Inst~Spectral E43/E50/E53/E54 exactly: 250 samples and \SI{0.375}{\second}
per single-sided scan. (The DEFAULT implementation is double-sided ---
500 samples, \SI{0.75}{\second} --- for complex-visibility recovery; see the
flag below.) The roof-mirror stage moves half the OPD
(mechanical travel $=\Delta/2$, [RM]~E51) --- \emph{not} the $\Delta/4$ of
FIInS's cat's-eye design, a heritage trap explicitly commented in the code.
\end{physicsbox}

The sky is simulated on the \emph{natural} FTS grid
$\sigma_k = k\,\delta\sigma$ (42 channels in band), so a noiseless
interferogram transforms back to the input spectrum with zero interpolation
error --- gridding artefacts can never masquerade as physics in the validation
loop.

\begin{flagbox}
\textbf{Finding (the most important one): the scan must be double-sided.}
The spreadsheet budgets a single-sided scan ($\Delta = 0 \to \Delta_{\max}$,
250 samples). That achieves $R=100$ only for a \emph{real} spectrum --- the
implicit assumption of conventional FTS, where the source contributes no
phase. A double-Fourier interferometer measures \emph{complex} visibilities:
the fringe phase carries the source positions, and is precisely the science.
Information counting: 250 real samples cannot determine 42 channels
$\times$ 2 quadratures of independent complex visibility plus the DC level
--- over a one-sided window the cosine and sine components of neighbouring
channels are nearly degenerate. The simulator demonstrated this numerically:
with the single-sided scan, noiseless visibility recovery fails at the
10--100\% level at many uv points; with the symmetric scan
($-\Delta_{\max} \to +\Delta_{\max}$, 500 samples, \SI{0.75}{\second}),
recovery is exact to $2\times 10^{-12}$. Consequence: the per-uv-point dwell
time doubles (full uv pass $\sim$\SI{54}{\second} at the spreadsheet's spin
strategy), or resolution halves to $R=50$. This is a real scan-budget
correction to take back to the design.
\end{flagbox}

% ----------------------------------------------------------------------
\subsection{\file{sky.py} --- the scene}

A scene is a list of compact sources; each source has a position, an angular
size, and a spectrum (blackbody continuum normalised to a flux density in Jy
at a reference wavelength, plus an optional Gaussian emission line for
validating the spectroscopy axis).

\begin{decisionbox}
\textbf{Discrete sources instead of a pixel cube.} PyFIInS generated an
$n_\mathrm{pix}\times n_\mathrm{pix}\times n_\sigma$ sky cube and extracted
visibilities by FFT. For a two-element instrument doing source separation,
the van Cittert--Zernike integral collapses to a finite sum over sources
(next subsection) which is evaluated \emph{exactly} --- faster, and free of
pixelisation error. The interface (\code{visibility(uv, sigma)}) is
unchanged, so a pixel-grid scene loader for extended emission can be added
later without touching the rest of the pipeline.
\end{decisionbox}

The convenience constructor \code{binary\_scene()} builds the canonical test:
an unequal binary (default $5\times10^4$ and $2\times10^4$~Jy, separation
8~arcsec) with an emission line on the fainter component --- the double-slit
configuration of Grainger Eqs.~11--13, with the line testing that spectra are
recovered \emph{per source}, which is the entire point of double-Fourier.

% ----------------------------------------------------------------------
\subsection{\file{beams.py} --- the Airy beam and the parity envelope}

\begin{physicsbox}
Each collector is a uniformly illuminated circular aperture, so its
amplitude beam is the jinc pattern $b(\theta) = 2J_1(x)/x$ with
$x = \pi D \sin\theta/\lambda$, first null at $1.22\lambda/D$
($\approx 50$~arcsec at \SI{10}{\micro\metre}) --- PyFIInS's `Perfect' beam
option; GRASP/measured beams can be slotted in behind the same interface.

\textbf{Parity-mismatch envelope (new calculation).} If the two arms have
opposite reflection parity, arm~2 delivers a mirror-flipped image: a source
at $+\theta_x$ in arm~1 arrives where $-\theta_x$ arrives in arm~2, so the
two single-arm PSFs for the same sky point are displaced by $2\theta_x$ and
the fringe contrast is their overlap. By Parseval, the overlap of two
displaced PSFs is the Fourier transform of the pupil intensity (a uniform
disc) at displacement $2\theta_x$:
\begin{equation}
\mu(\theta_x) = \frac{2J_1(z)}{z}, \qquad
z = \frac{2\pi D\,\theta_x}{\lambda},
\end{equation}
unity on axis and zero at $\theta_x = 0.61\lambda/D$ --- \emph{half} the Airy
radius. A parity mismatch therefore halves the usable field of view in the
flip direction while costing nothing on axis. The envelope multiplies the
visibility integrand when \code{parity\_mismatch=True}, making the
field-of-view penalty directly simulatable --- the quantitative version of
the extra-mirror question.
\end{physicsbox}

% ----------------------------------------------------------------------
\subsection{\file{uvcoverage.py} --- the baseline schedule}

Implements the spreadsheet's uv strategy ([RM]~Spacecraft rows 22--31):
spin rate $\omega(b) = (d/2)/(t_\mathrm{scan}\,b/2)$, i.e.\ one interferogram
scan per dish-diameter of arc motion, which makes the angular step between uv
samples $\delta\phi = d/b$ --- exactly the uv-plane Nyquist criterion (adjacent
samples separated by one aperture diameter). One half-circle per baseline
length suffices because the sky is real ($V(-\mathbf u) = V^*(\mathbf u)$);
the baseline is then stepped by $2d$ ([RM]'s choice: each dish moves one
diameter per half-turn; this undersamples radially by $2\times$ relative to
Nyquist --- flagged, parameterised). The frozen-baseline approximation
(baseline fixed during each scan) is consistent with the design's own
smearing tolerance; intra-scan fringe smearing is a documented missing
systematic.

% ----------------------------------------------------------------------
\subsection{\file{radiometry.py} --- background, noise, sensitivity}

\begin{physicsbox}
Each warm emitter contributes
$P = \int \epsilon_\mathrm{eq}(\sigma)\, B_\sigma(T)\, A\Omega(\sigma)\,
\eta_\mathrm{downstream}\, \eta_\mathrm{det}\, d\sigma$; the photon NEP is
the Poisson expression $\mathrm{NEP}^2 = \int 2 E_\mathrm{ph}\,
P_\mathrm{abs,\sigma}\, d\sigma$ (Bose excess $<1\%$ at
$h\nu/kT \approx 6$); total NEP adds the detector dark NEP in quadrature;
the rms noise on one OPD sample of duration $t_s$ is
$\mathrm{NEP}\sqrt{1/(2t_s)}$. The module reproduces the spreadsheet's
three background loadings to ratio $1.00$ (environment
\SI{5.38e-9}{\watt}, telescope \SI{2.52e-10}{\watt}, cold optics
\SI{1.02e-9}{\watt}), which validates the emissivity chain end to end.
\end{physicsbox}

\begin{flagbox}
\textbf{The spreadsheet's photon-NEP rows are internally inconsistent with
its own loadings.} Applying $\mathrm{NEP}=\sqrt{2PE_\mathrm{ph}}$ to its SUN
loading (\SI{5.4e-9}{\watt}) gives $\sim$\SI{5e-15}{\watt\per\sqrt{\hertz}},
yet its NEP cell says \SI{6.5e-17}{} --- a factor $\sim$100, i.e.\ a factor
$10^4$ in power, exactly the $\mathrm{cm}^2/\mathrm{m}^2$ ratio. The
`Wavelengths' sheet that feeds those cells carries flux in
$[\,1/\mathrm{cm}^2/\mathrm{sr}\,]$ units against etendue in
$\mathrm{m}^2\,\mathrm{sr}$. Our first-principles photon NEP is
\SI{1.1e-14}{\watt\per\sqrt{\hertz}}, not \SI{9.2e-17}{}. To check with GS;
the spreadsheet's own loading rows (which we reproduce exactly) support the
larger number.
\end{flagbox}

\begin{flagbox}
\textbf{Detector NEP is the single most important missing number.} The
spreadsheet's NEP(Dark) cells are empty. The placeholder
(\SI{5e-14}{\watt\per\sqrt{\hertz}}, a representative TEC-cooled mid-IR
photodetector, $D^*\!\sim\!10^{11}$\,cm\,$\sqrt{\mathrm{Hz}}$/W, 50-\si{\micro\metre}
pixel) makes the system \emph{detector-noise limited} ($\sim$5$\times$ above
even our corrected photon NEP, $\sim$500$\times$ above the spreadsheet's
figure). At that noise the 5$\sigma$ per-channel sensitivity is
$\sim 9\times10^3$~Jy per full uv pass (worst band): only the very brightest mid-IR
sources (Betelgeuse-class and brighter) are accessible. A cryogenic detector
(NEP $\sim 10^{-16}$) changes the conclusion completely. This is the gate
question for the noise budget.
\end{flagbox}

% ----------------------------------------------------------------------
\subsection{\file{forward.py} --- the double-Fourier forward model}

\begin{physicsbox}
The detector power at OPD $\Delta$ and baseline $\mathbf b$ is the
double-Fourier measurement equation (Grainger Eq.~15, per channel):
\begin{equation}
P(\Delta) = \sum_\sigma \delta\sigma\, A\,\eta(\sigma)\!\left[
\tfrac{t_1+t_2}{2}\,S_0(\sigma) + \sqrt{t_1 t_2}\,
\mathrm{Re}\!\left\{V(\sigma)\, e^{\,2\pi i \sigma\Delta}\right\}\right]
+ P_\mathrm{bg},
\end{equation}
with the visibility given by the van Cittert--Zernike sum over the discrete
sources,
\begin{equation}
V(\sigma) = \sum_k S_k(\sigma)\, b^2(\theta_k,\sigma)\,
\mu(\theta_{k,x})\, \tau_k(\sigma,b)\,
e^{-2\pi i \sigma\,(\mathbf b\cdot\boldsymbol\theta_k)},
\end{equation}
where $b^2$ is the two-collector beam product, $\mu$ the parity envelope,
$\tau_k$ a Gaussian resolved-size taper, and the DC term $S_0$ uses the
intensity beam. The two Fourier kernels --- spatial
$e^{-2\pi i\sigma\,\mathbf b\cdot\boldsymbol\theta}$ and spectral
$e^{2\pi i\sigma\Delta}$ --- share the same $\sigma$ inside one expression:
that is, literally, the ``double Fourier''.

The detector's \SI{500}{\micro\second} time constant low-pass filters the
fringes: during a scan, channel $\sigma$ appears at temporal frequency
$f = \sigma\,\dot\Delta$ (each colour becomes its own audio tone) and is
attenuated by the single-pole MTF $1/\sqrt{1+(2\pi f\tau)^2}$ ---
0.88--0.97 across the band at the design scan speed. Included in the forward
model; divided back out in reduction as a known calibration.
\end{physicsbox}

Noise is injected per sample at the radiometric level; \code{seed=None}
disables it for architecture validation. Pointing jitter has an explicit
hook but no model yet --- PyFIInS used measured \emph{Herschel} jitter
spectra, which have no relevance to a reaction-wheel CubeSat; a CubeSat ADCS
jitter spectrum is a placeholder input.

% ----------------------------------------------------------------------
\subsection{\file{reduction.py} --- interferogram to visibility spectrum}

The first of the two inverse Fourier transforms.

\begin{decisionbox}
\textbf{Least squares instead of FFT.} PyFIInS circularly shifted each
interferogram so zero OPD sat at index 0, then FFT'd (with a code comment
``must ask David N how to do this properly''). Here the OPD axis is known
exactly, so each scan is solved in one step as a linear model
\begin{equation}
P(\Delta_j) = c_0 + \sum_k \big[a_k \cos(2\pi\sigma_k\Delta_j)
 + b_k \sin(2\pi\sigma_k\Delta_j)\big],
\qquad \hat V_k = \frac{a_k - i\,b_k}{g_k},
\end{equation}
with the known gain $g_k$ (area $\times$ efficiency $\times$
$\sqrt{t_1t_2}$ $\times$ $\delta\sigma$ $\times$ detector MTF) divided out so
$\hat V$ returns in physical flux-density units. Advantages over the FFT
route: no zero-OPD shifting or phase-correction step at all; the constant
term (DC + background) is fitted simultaneously rather than subtracted (a
mean-subtraction biases channels whose cosines have a half-integer number of
cycles in the window); and the method generalises unchanged to irregular OPD
sampling when delay-line velocity errors are added later. Cost: one
$500\times 85$ \code{lstsq} per scan --- negligible.
\end{decisionbox}

% ----------------------------------------------------------------------
\subsection{\file{imaging.py} --- the second transform and source separation}

Two products. First, the \textbf{dirty map} (diagnostic only):
$I_d(\boldsymbol\theta) = \langle \mathrm{Re}\{\hat V\,
e^{+2\pi i\sigma\,\mathbf b\cdot\boldsymbol\theta}\}\rangle_{\mathrm{uv},\sigma}$,
plus the matching dirty beam. With $\sim$50--100 uv points on 2--3 rings the
dirty beam has strong sidelobes; the map serves to find starting positions,
not to do photometry. (Hogbom CLEAN, which PyFIInS carried but had commented
out of its own pipeline, is deliberately omitted: for a few-source scene,
parametric fitting is simpler and statistically better.)

Second, the \textbf{science pipeline} --- source separation + spectroscopy:
\begin{enumerate}
\item initial positions from the band-averaged dirty map peaks (grid search
with parabolic sub-pixel refinement);
\item position refinement by separable least squares (variable projection):
at fixed positions the per-channel model
$V(u,v) = \sum_k S_k\, e^{-2\pi i\sigma\,\mathbf b\cdot\boldsymbol\theta_k}$
is \emph{linear} in the $S_k$, so the best-fit fluxes and residual are cheap;
Nelder--Mead then minimises over the $2n_\mathrm{src}$ position coordinates.
Dirty-map peaks alone are accurate to a fraction of a pixel, which leaks
phase errors into the spectra ($\sim$8\% at 0.3~arcsec); the refinement
closes the noiseless loop to machine precision;
\item per-channel linear solve at the refined positions, then division by
the known primary beam (and parity envelope) at each position. This inverts
the double-slit cosine of Grainger Eqs.~11--13 for amplitude \emph{and}
phase, rather than just fitting its period.
\end{enumerate}

% ----------------------------------------------------------------------
\subsection{\file{validate.py} and \file{run\_simulation.py}}

\code{validate} computes the loop metrics (position errors in arcsec;
fractional rms between recovered and true spectra, matched source by source)
and writes the diagnostic figures: uv coverage, noiseless and noisy
interferograms (the white-light fringe at $\Delta=0$ is clearly visible),
dirty map + dirty beam with true and recovered positions overplotted, and
the input-vs-recovered spectra. \code{run\_simulation.py} chains everything
and prints the radiometric cross-check against the spreadsheet, the
sensitivity summary, and the metrics for both the noiseless and noisy runs.

% ======================================================================
\section{Results of the reference run}
% ======================================================================

This reference run exercises the \emph{future} multi-baseline
rotation+contraction pipeline, not the confirmed mission: the confirmed
single-baseline design (single fixed ${\sim}\SI{0.4}{\metre}$, delay-line
scan only, no uv rotation or contraction) is characterised in the
\file{run\_crater.py} worked example below.

Default parameters; unequal binary ($5\times10^4$ / $2\times10^4$~Jy,
8~arcsec separation, PA $30^\circ$, emission line on the fainter source);
72 uv points on 3 baseline rings (0.5, 0.398, \SI{0.3}{\metre}),
\SI{54}{\second} of observation; 5$\sigma$ per-channel sensitivity
$9\times10^3$~Jy for the full pass, worst band (after the combiner-gain and detector-MTF fixes).

\begin{center}\small
\begin{tabular}{@{}lll@{}}
\toprule
\textbf{Quantity} & \textbf{Noiseless} & \textbf{Noisy (placeholder NEP)}\\
\midrule
position error, source A & $<0.001''$ & $0.06''$\\
position error, source B & $<0.001''$ & $0.10''$\\
spectrum rms, source A & $<10^{-6}$ & 8.6\%\\
spectrum rms, source B & $<10^{-6}$ & 17.5\%\\
\bottomrule
\end{tabular}
\end{center}

The noiseless loop closing exactly is the architecture validation. The noisy
numbers say something real: even at $5\times10^4$~Jy (brightest-star
territory) the per-channel SNR with a TEC-class detector is modest --- the
detector NEP question (Flag~2) is where the science capability will be
decided.

% ======================================================================
\section{Single-baseline analysis and the three target cases}
% ======================================================================

The full pipeline above rotates the baseline to fill the uv plane. But the
\emph{building block} of the whole technique is what happens at ONE baseline:
hold the telescopes fixed, scan the delay line once, and look at what comes
out. This is the right level at which to choose a science target and ask
``can we make out the feature?'', and it is implemented separately in
\file{singlebaseline.py} with three ready-to-run targets in \file{scenes.py}
(\code{python run\_cases.py [name]}). The cases are independent --- pick one
by name and run it without touching the others.

\begin{physicsbox}[title={\textsc{Peak, ripples, convolution: the FTS picture
at one baseline}}]
A single delay-line scan produces an \textbf{interferogram} --- detector power
vs OPD: a central PEAK (the white-light fringe, where every colour adds in
phase at zero OPD) with FRINGES/RIPPLES around it. The ripples are the
source's spectrum and spatial structure imprinted on the delay axis.
Fourier-transforming the interferogram recovers the spectrum (the
``peak''). But the scan is \emph{finite} (length OPD$_{\max}$), which is the
same as multiplying the true interferogram by a boxcar window --- and a
boxcar in OPD is a SINC in wavenumber. So every true spectral peak comes back
\textbf{convolved with a sinc that has ripples (sidelobes at $\sim$13\%)}:
the instrument line shape. That is the ``ripples in the convolution''. The
ratio (scan length)/(feature scale) sets how many ripples and how sharp the
recovered peak is --- the ``10\,m feature, 100\,m measuring scale'' picture,
on whichever axis (a longer scan resolves a narrower line; a longer baseline
resolves a smaller angle). \textbf{Apodization} (tapering the window, here
Hann) trades the ripples for a broader peak; both curves are shown.
\end{physicsbox}

To make the line-shape ripples visible at all, the source is simulated on a
spectral grid finer than the FTS channels (\code{Case.spectral\_oversample})
and the band edge is given a smooth roll-off (the real band-defining filter,
\code{singlebaseline.band\_filter}) instead of a razor edge that would ring
spuriously. This deliberately exercises the one class of FTS systematics that
the on-channel-grid validation loop cannot (it was noted there as a
structural blind spot).

\paragraph{The three cases (each a candidate target and a physics regime).}

\begin{center}\small
\begin{tabular}{@{}p{2.6cm}p{4.6cm}p{5.6cm}@{}}
\toprule
\textbf{Case} & \textbf{Candidate target} & \textbf{What the single baseline shows}\\
\midrule
\code{compact\_line} & a star, or an atmospheric gas line (an unresolved
point with a spectral feature) & the SPECTRAL axis: the line is recovered at
1003~cm$^{-1}$ as a sinc (the instrument line shape) with ripples;
apodization smooths them. Visibility is flat (unresolved).\\[4pt]
\code{binary} & two craters, a binary star, two thermal hotspots & SOURCE
SEPARATION: the 24$''$ separation imprints a cosine RIPPLE across the
spectrum ($\sim$2 cycles), period $1/(b\,\theta_{\rm sep})$; reading it gives
the separation. This is the double-slit result on one baseline.\\[4pt]
\code{resolved\_disk} & a planetary/lunar disk or limb, a wide thermal
region & RESOLUTION: the visibility falls ($0.56\!\to\!0.27$ across the band),
more at bluer wavenumbers --- the source is resolved out, and the recovered
spectrum is tilted/suppressed accordingly.\\
\bottomrule
\end{tabular}
\end{center}

Each case writes a three-panel figure (interferogram / recovered spectrum vs
truth / visibility) to \code{outputs/cases/}. Together they answer the
professor's question directly: \code{compact\_line} shows whether a spectral
feature survives the line shape; \code{binary} whether two features are
separable and how far apart; \code{resolved\_disk} whether an extended target
is resolved and at which colours.

\begin{flagbox}[title={\textsc{What one baseline cannot do}}]
At a single baseline the recovered spectrum is $S(\sigma)\,V(\sigma)$ --- the
true spectrum times the spatial visibility at that one baseline --- and the
two CANNOT be separated from one baseline alone. A cosine ripple across the
band (the \code{binary} case) could equally be a genuine spectral modulation;
a tilt (\code{resolved\_disk}) could be a redder source. Disentangling
spatial from spectral is exactly what baseline diversity --- the full uv
rotation of the main pipeline --- is for. The single-baseline analysis is the
per-baseline building block; the full pipeline stacks many of them, and only
then can position and per-source spectrum be told apart (the
``$<10^{-6}$ noiseless'' result of the previous section).
\end{flagbox}

% ======================================================================
\section{Worked example: single-baseline lunar-crater detection}
% ======================================================================

The CONFIRMED instrument as scoped (GS, June 2026) is a \emph{single fixed
baseline, one axis, no baseline movement} (the delay line still scans for the
spectrum; only the baseline vector is held). That rules out imaging --- there
is no uv coverage to fill --- but it supports \emph{detection and one-axis
sizing}, the regime of the ship-on-sea example. (The multi-baseline
rotation+contraction pipeline in \code{run\_simulation.py} can image, but is a
future/non-mission configuration, not the confirmed design.) This section runs that case for a lunar crater
and uses it as the evidence for the rigid-bench-vs-boom decision (open
question~1).

\subsection{\file{extended.py} --- pixel-map scenes}

\begin{decisionbox}
A crater on the Moon is an \emph{extended brightness map}, not a few compact
sources, so \file{extended.py} adds a new scene type: a temperature map
$T(x,y)$ on an angular grid, each pixel a blackbody. Its visibility is the
van Cittert--Zernike sum over pixels,
\[
V(\sigma) = \sum_{\mathrm{pix}} \epsilon\, B_\sigma(T_{\mathrm{pix}})\,
b^2(\theta_{\mathrm{pix}})\,\mu(\theta_{\mathrm{pix}})\,
e^{-2\pi i\sigma(\mathbf b\cdot\boldsymbol\theta_{\mathrm{pix}})}\,\mathrm{d}\Omega,
\]
the same sum \file{forward.py} uses, with sources replaced by pixels of solid
angle $\mathrm{d}\Omega$. To keep one source of truth, the interferogram
formula (gains, MTF, ZPD, noise, saturation) was refactored into a single
shared function, \code{forward.interferogram\_from\_visibility}, that both
the discrete and extended paths call --- so any combiner-normalisation fix
(the combiner factor-of-2, now fixed in \code{forward.fringe\_gain})
propagates to both. The loader is verified
against the discrete forward model in the point-source limit (a single hot
pixel returns $|V|/S_0 = 1$ with the exact geometric phase ramp), and the
noiseless interferogram round-trips to the analytic visibility at $10^{-6}$.
\end{decisionbox}

\subsection{The detection physics}

\begin{physicsbox}
A uniform surface patch filling the beam washes the fringes out: its power
sits in the un-fringing DC term and $V\to0$. A crater breaks that uniformity,
lifting $|V|$ above the uniform-patch floor --- that lift, relative to the
radiometric noise, \emph{is} the detection. The crater's diameter along the
baseline is read from the \emph{shape} of $|V|(\sigma)$: because $u=b/\lambda$
sweeps a radial line across the band, one fixed baseline samples part of the
visibility's sinc fall-off. The fixed baseline is therefore a matched filter
--- a crater near $\lambda/b$ gives the strongest, most size-diagnostic
signal; one far smaller is unresolved (a small flat $V$); one filling the
beam washes out again.
\end{physicsbox}

\subsection{Result --- the baseline decision}

A \SI{30}{\kilo\metre} crater (cool floor) in a uniform sunlit patch, observed
at three fixed baselines (\code{run\_crater.py}). The single-dish field of
view is $\sim50''\approx\SI{92}{\kilo\metre}$ at the Moon. Because the Moon is
extremely bright, every case is detected with SNR $\gg5$ (after the combiner
fix, $\sim$20--400) --- the study is \emph{resolution}-limited, not
noise-limited, so the detector-NEP placeholder does not change the
conclusion (and the combiner normalisation, now corrected, only rescales the
absolute flux). The differentiator is \emph{sizing}, which needs the crater
larger than the resolution element:

\begin{center}\small
\begin{tabular}{@{}lll@{}}
\toprule
\textbf{Baseline} & \textbf{Resolution} & \textbf{Can size\ldots}\\
\midrule
\SI{0.125}{\metre} (rigid bench) & $20''$ ($\sim$\SI{37}{\kilo\metre}) &
only $\sim$\SI{90}{\kilo\metre} (Copernicus-class); 10--30 km detected, not sized\\
\SI{0.25}{\metre} & $10''$ ($\sim$\SI{19}{\kilo\metre}) & $\sim$\SI{30}{\kilo\metre} and up\\
\SI{0.5}{\metre} (boom) & $5''$ ($\sim$\SI{9}{\kilo\metre}) & down to $\sim$\SI{10}{\kilo\metre}\\
\bottomrule
\end{tabular}
\end{center}

So the rigid-vs-boom choice reduces to a science question: \emph{do we need to
measure crater sizes (boom), or only detect that structure is present (the
rigid bench suffices)?} A nuance the figures show: the rigid bench gives a
\emph{stronger} detection signal ($|V|/S_0\approx0.025$ vs $0.005$ at
\SI{0.5}{\metre}), because a longer baseline resolves --- and so washes out
--- more of the crater. Short baseline: higher contrast, coarse sizing; long
baseline: finer sizing, lower contrast (still trivially detected here).

Caveat: the crater is modelled as a uniform cool disk (\SI{390}{\kelvin}
surface, \SI{250}{\kelvin} floor); a real crater has rim/floor/central-peak
structure. The technique detects any departure from uniformity, so the
conclusion holds, but the exact signature would differ. Figures live in
\code{outputs/crater/} (the patches, the $|V|(\sigma)$ signature per baseline,
and the SNR-vs-size-vs-baseline grid).

% ======================================================================
\section{All flags and open questions, collected}
% ======================================================================

\begin{enumerate}
\item \textbf{Baseline geometry (RESOLVED).} The confirmed design is a
SINGLE FIXED baseline of ${\sim}\SI{0.4}{\metre}$ (detection + one-axis
sizing, not imaging). The spreadsheet's \SIrange{0.3}{0.5}{\metre} with
contraction (deployable boom / formation flight + contraction mechanism)
is a future multi-baseline option; the \SI{125}{\milli\metre} rigid bench is
the detection-only alternative. Angular resolution spans $4\times$
($\lambda/b$ at \SI{10}{\micro\metre}: $16.5''$ at \SI{0.125}{\metre} vs
$4.1''$ at \SI{0.5}{\metre}). The \code{Geometry} defaults keep the
spreadsheet \SIrange{0.3}{0.5}{\metre} only for the legacy
\code{uvcoverage.py} schedule; set
\code{baseline\_min\_m = baseline\_max\_m} to a single value (0.4 design,
0.125 rigid bench) for single-baseline runs.
\item \textbf{Detector NEP} ([PH], the gate number): empty in the
spreadsheet; placeholder \SI{5e-14}{\watt\per\sqrt{\hertz}} makes the system
detector-noise limited and sets the $\sim 9\times10^3$~Jy sensitivity floor.
\item \textbf{Arm parity}: matched or mismatched handedness after the extra
fold? Mismatch halves the usable field ($\mu$ envelope); throughput effect
is negligible ($4\times10^{-5}$). Needs a reflection count per arm from the
optical prescription.
\item \textbf{Scan budget (RESOLVED: double-sided is the default).}
Single-sided scanning cannot recover complex visibilities at $R=100$; the
default is therefore double-sided, doubling the dwell time
(\SI{0.75}{\second}/scan) --- found numerically, see \file{fts.py} flag.
\item \textbf{Spreadsheet photon-NEP inconsistency}: loadings and NEPs
disagree by $\sim 100\times$ (suspected cm$^2$/m$^2$ slip); our
first-principles value is \SI{1.1e-14}{\watt\per\sqrt{\hertz}}.
\item \textbf{Stale labels in the spreadsheet}: the pixel-FoV / etendue
cells are labelled ``@ 8\,\si{\micro\metre}'' but their cached values
correspond to \SI{10}{\micro\metre}; the band itself (8--12~\si{\micro\metre},
mid-IR rather than far-IR) is worth confirming as intentional. The
spreadsheet's source is currently an \SI{1800}{\kelvin} lab circuit against
a \SI{293}{\kelvin} scene (testbed configuration), and its mass/power tables
are inherited from the FISICA strawman (1.6~m telescopes, 14~m booms) ---
neither describes the CubeSat.
\end{enumerate}

% ======================================================================
\section{What is deliberately not modelled (yet)}
% ======================================================================

Kept out to stay at the right level of sophistication, each with its natural
insertion point:

\begin{itemize}
\item \textbf{Pointing jitter} --- hook in \code{observe\_scan}; needs a
CubeSat ADCS (reaction-wheel) jitter spectrum, not the Herschel spectra
PyFIInS used.
\item \textbf{Delay-line velocity errors} --- enter as an irregular OPD
axis; the least-squares reduction already accepts arbitrary OPD sampling.
\item \textbf{Intra-scan baseline rotation (fringe smearing)} --- bounded by
the design's own spin criterion; a smearing factor per channel would slot
into \code{forward.detector\_mtf}'s position in the chain.
\item \textbf{Extended emission} --- a pixel-grid scene class behind the
same \code{visibility()} interface. Related, verified numerically: the
point-source extractor under-reports a \emph{resolved} source's flux by
the uv-averaged size taper (a 3-arcsec FWHM source on the default
baselines returns $0.26\times$ its true flux); fitting a size parameter
per source is the upgrade path. Channels where the beam/parity
correction would divide by $\sim$zero are returned as NaN rather than
fabricated.
\item \textbf{Measured beams} --- GRASP import behind
\code{beams.amplitude\_beam}.
\item \textbf{Thermal drift of the ZPD} --- the static residual offset is
now modelled (\code{FTSScan.zpd\_offset\_cm}); its slow per-scan \emph{drift}
is not, and would enter as a time-dependent offset in the same place.
\item \textbf{Beam walk on the roof mirror} --- the \SI{0.5}{inch}
collimated beam translates across the moving roof pair during a scan,
modulating coupling and adding a scan-position-dependent throughput; needs
the OAP prescription from the layout to quantify.
\item \textbf{Detector nonlinearity and clipping} --- the saturation hook
warns when a sample exceeds \code{saturation\_power\_w} ([PH]) but the
response is otherwise perfectly linear; a measured response curve would
slot in where the warning fires.
\end{itemize}

% ======================================================================
\section{Relation to the PyFIInS heritage}
% ======================================================================

\begin{center}\small
\begin{tabular}{@{}p{4.2cm}p{4.4cm}p{4.6cm}@{}}
\toprule
\textbf{Function} & \textbf{PyFIInS (FISICA)} & \textbf{cubesat\_sim}\\
\midrule
language & Python 2, Parallel Python & Python 3, numpy/scipy only\\
parameters & Excel workbooks parsed at runtime & source-traced dataclasses\\
sky & pixel cube + FFT & exact discrete-source VCZ sum\\
uv coverage & boom spirals (const-$L$, spiro) & spin + contraction rings\\
scan & double-sided, $\Delta/4$ cat's-eye & double-sided, $\Delta/2$ roof pair\\
errors & Herschel jitter spectra & radiometric noise; CubeSat jitter [PH]\\
spectra from IGs & shift + FFT & explicit-OPD least squares\\
image product & dirty cube (+ unused CLEAN) & dirty map + parametric
source separation\\
noise model & none in forward chain & background + photon + detector NEP\\
\bottomrule
\end{tabular}
\end{center}

The thesis-Chapter-4 module map survives intact --- sky, beams, uv,
FTS, observe, reduce, image, validate --- which is deliberate: anyone who
knows FIInS can navigate this code immediately, and anything learned here
maps back onto the heritage architecture.

\end{document}
